\subsection{Análise das Regras de Negócio}

A análise das regras de negócio concentrou-se nos critérios de classificação dos especialistas, nas alterações cadastrais e nas permissões associadas às operações de solicitação, aprovação e rejeição. O levantamento dessas regras permitiu estabelecer os requisitos comportamentais do novo módulo, considerando as restrições existentes e os possíveis impactos das operações sobre os registros do sistema.

Um dos principais desafios identificados decorre da organização do modelo de dados da aplicação: a categoria do especialista é definida no modelo \textit{User}, e não diretamente no modelo \textit{Specialist}. Embora os especialistas compartilhem uma estrutura de usuário comum, suas categorias determinam especializações distintas, com atributos específicos para cada tipo. Essa organização exige que as operações de transição considerem as relações entre os modelos envolvidos, de modo a preservar a integridade dos dados durante a alteração cadastral.

Outro aspecto relevante consiste em assegurar a rastreabilidade das alterações e a possibilidade de preservar o estado anterior dos registros durante o fluxo de solicitação. Para isso, é necessário distinguir o registro de uma solicitação da efetivação da mudança, condicionando esta última ao resultado do processo de análise. Adicionalmente, as transições devem considerar os critérios de classificação e pontuação aplicáveis a cada categoria, a fim de evitar inconsistências decorrentes da alteração dos dados cadastrais.